\subsection{The bare prime budget for every fixed cyclotomic type}

\begin{theorem}[DP64: fixed cyclotomic types and saturated integer budgets]\label{DP64}
Fix a CM subfield $E\subset\mathbf Q(\zeta_m)$, its embedding set
$\Omega$, a CM type $\Phi\subset\Omega$, and a finite prime set $S$
containing the divisors of $m$.
For $p\notin S$, let $\sigma_p\in\operatorname{Gal}(E/\mathbf Q)$
be the restriction of $\zeta_m\mapsto\zeta_m^p$.
For $i\in\Phi$ set
\[
 c_{i,p}=\sum_{r\ge1}\frac{\mathbf1_{\sigma_p^ri\notin\Phi}}{p^r},
 \qquad E_{i,p}(n)=\sum_{r\ge1}\mathbf1_{\sigma_p^ri\notin\Phi}
                     \left\lfloor\frac n{p^r}\right\rfloor.
\]
These are the DP40/43 bare directional slopes and integer lower
certificates for the associated integral cyclic displays; applying
them to geometry requires those displays and fixed global primitive
eigendirections.
For arbitrary nonnegative integers $n_i$, held fixed as $p$ varies,
put $N=\sum_i n_i$. Uniformly in the type, labels, orders and $X\ge2$,
\begin{align}
 \sum_{\substack{p\le X\\p\notin S}}\log p\sum_i n_i c_{i,p}
 &=\frac N2\log X+O_{m,S}(N),\label{eq:generalcmlinear}\\
 \sum_{\substack{p\le X\\p\notin S}}\log p\sum_i E_{i,p}(n_i)
 &=\frac12\sum_i n_i\log\bigl(\max\{1,\min(X,n_i)\}\bigr)
       +O_{m,S}(N).\label{eq:generalcmfloors}
\end{align}
Thus data in the same fixed cyclotomic field with the same order
multiset have budget difference $O_{m,S}(N)$.
If at most $d$ directions are used, even arbitrary distributions of
the same total order satisfy
\begin{equation}\label{eq:generalcmentropy}
 \sum_{\substack{p\le X\\p\notin S}}\log p\sum_i E_{i,p}(n_i)
 =\frac N2\log\bigl(\max\{1,\min(X,N)\}\bigr)
       +O_{m,S,d}(N).
\end{equation}
The constants are not asserted uniform in growing $m,S,d$.
The orders may vary with an outer parameter, but are fixed during
each prime sum.
\end{theorem}

\begin{proof}
The known analytic input is Rudnick, \emph{Primes in Arithmetic
Progressions: Fixed Modulus}, May 26, 2015, \S2, printed p.~2,
equation~(1): for fixed $m$ and $(a,m)=1$,
\[
 \sum_{\substack{p\le Y\\p\equiv a\pmod m}}\frac{\log p}{p}
 =\frac{\log Y}{\varphi(m)}+O_m(1).
\]
The quantitative fixed-modulus PNT immediately above it also gives
$\sum_{p\le Y}\log p=O(Y)$.
These are cited known results, independent of RH.\footnote{The exact
primary page was reopened and read:
\url{https://www.math.tau.ac.il/~rudnick/courses/sieves2015/fixed\%20progressions.pdf}.}

Put $b_i(a)=\mathbf1_{\sigma_a i\notin\Phi}$.
Restriction from $(\mathbf Z/m)^\times$ onto
$G=\operatorname{Gal}(E/\mathbf Q)$ is surjective with equal fibres,
and $G\to\Omega$, $\sigma\mapsto\sigma i$, is bijective.
Exactly half of $\Omega$ lies outside $\Phi$, so
\[
 \frac1{\varphi(m)}\sum_{a\in(\mathbf Z/m)^\times}b_i(a)=\frac12.
\]
Summing the analytic input in the finitely many classes, and deleting
$S$, gives
\begin{equation}\label{eq:generalcmfirst}
 \sum_{\substack{p\le Y\\p\notin S}}b_i(p)\frac{\log p}{p}
 =\frac12\log Y+O_{m,S}(1).
\end{equation}
The sum of the absolute class errors bounds the error for every
$0$--$1$ function $b_i$, giving the stated uniformity in types and
labels. Abelianness makes $\sigma_p$ a well-defined element specified
by the actual residue class; there is no conjugacy-class sampling
assumption hidden in this step.

The higher-power tail satisfies
\[
 0\le c_{i,p}-\frac{b_i(p)}p\le\frac1{p(p-1)}.
\]
The series $\sum_p\log p/[p(p-1)]$ converges by comparison with
$\sum_{n\ge2}\log n/[n(n-1)]$.
Multiply \eqref{eq:generalcmfirst} by $n_i$, add the tails and sum
over $i$ to prove \eqref{eq:generalcmlinear}.
Entirely selected orbits contribute zero, as their formulas require.

For the integer assertion fix $i$ and write $n=n_i$.
If $n=0$ or $1$, every floor and the asserted main term vanish.
For $n\ge2$, put $Y=\min(X,n)$.
Replacing each first-power floor by $n/p$ makes an error bounded
by $\sum_{p\le Y}\log p=O(Y)=O(n)$; therefore
\[
 \sum_{\substack{p\le Y\\p\notin S}}
 b_i(p)\left\lfloor\frac np\right\rfloor\log p
 =\frac n2\log Y+O_{m,S}(n).
\]
The full higher-power contribution is nonnegative and at most
\[
 n\sum_p\log p\sum_{r\ge2}p^{-r}=O(n).
\]
Adding proves \eqref{eq:generalcmfloors}.
Its main term saturates at $X=n$ for each direction, matching the
finite floor sum; no indefinitely growing fixed-order main term is
asserted. Subtracting the formula for a common order multiset proves
the first difference statement.

Finally assume $N>0$, let $I=\{i:n_i>0\}$, put $w_i=n_i/N$,
and set $L=\min(X,N)\ge1$.
For every $i\in I$,
\[
 0\le\log L-\log\min(X,n_i)\le\log(N/n_i).
\]
Indeed, if $X\le n_i$ the difference is zero; otherwise it is
$\log(L/n_i)$ with $L\le N$.
Multiplying by $n_i$ and adding yields
\[
 0\le N\log L-\sum_{i\in I}n_i\log\min(X,n_i)
 \le N\sum_{i\in I}w_i\log(1/w_i)\le N\log d.
\]
For the last inequality, concavity of $\log$ gives
$\sum_iw_i\log(1/w_i)\le\log\sum_iw_i/w_i=\log|I|$.
When $N=0$ every term is zero.
Together with \eqref{eq:generalcmfloors} this proves
\eqref{eq:generalcmentropy}.
\end{proof}

These comparisons concern bare lower certificates and do not upper-bound
selected theta values or the source's actual gain. Changes of order $N$
can still matter after leading terms cancel. Changing the conductor
requires control of constants not supplied here.

Prime-dependent choice is a distinct possibility. For
$\Phi=\{1,2,3\}\subset(\mathbf Z/7)^\times$,
$\max_{i\in\Phi}\mathbf1_{ai\notin\Phi}$ equals zero for $a=1$
and one for each other class: for $a=2,3,4,5,6$ the sets $a\Phi$ are
respectively $\{2,4,6\}$, $\{2,3,6\}$, $\{1,4,5\}$,
$\{1,3,5\}$ and $\{4,5,6\}$.
Its class mean is therefore $5/6$, exceeding $1/2$.
This finite calculation is not a global realization of varying local
directions with negligible cost. That construction and its norm and
selection constraints require a separate argument. Nonabelian and
reflex-field prime sampling are likewise outside the stated hypotheses.
